\section{Theory}
\label{sec:theory}

\begin{figure}[tb]
  \centering
  \includegraphics[width=\linewidth]{fig_method_overview.pdf}
  \caption{\textbf{\evq{} frequency allocation.} Geometric RoPE spaces channels
  uniformly in log-frequency; \evq{} applies a closed-form inverse-CDF warp that
  reallocates the same \emph{number} of channels toward wavelengths completing a
  training-window cycle. Rotation, architecture, and optimiser are unchanged.}
  \label{fig:method-overview}
\end{figure}

Figure~\ref{fig:method-overview} summarises the construction. We keep three
layers apart, because they carry different epistemic weight: \emph{exact}
statements about the geometry a finite table can express
(\S\ref{sec:subspaces}--\S\ref{sec:coordinate}); a \emph{modelling choice} that
turns one of them into a closed-form schedule (\S\ref{sec:construction}); and an
\emph{empirical} operating point. No experimental claim in
\S\ref{sec:experiments} rests on the first layer alone --- indeed
\S\ref{sec:coordinate} shows it cannot.

\subsection{A finite table is a set of two-dimensional subspaces}
\label{sec:subspaces}

Let $\phi \in [0,1]$ parameterise log-frequency, $\omega(\phi) = b^{-\phi}$, let
$\drot$ be the number of rotated dimensions and $K = \drot/2$ the number of
rotary pairs ($\drot{=}\dhd$ for standard MHA, smaller for compressed-RoPE
architectures such as MLA).

A rotary pair contributes $f_\omega(\Delta) = C\cos(\omega\Delta) +
D\sin(\omega\Delta)$ to the relative-position attention logit, with $C,D$ set by
content. The object attached to $\omega$ is therefore the two-dimensional
subspace
\begin{equation}
V_\omega=\operatorname{span}\{\cos(\omega\Delta),\ \sin(\omega\Delta)\},
\label{eq:subspace}
\end{equation}
not a cosine feature: an analysis that fixes one content phase sees a slice of
$V_\omega$ and is not invariant to rotation inside the pair. Two channels are
redundant to the extent that their subspaces are. With
$x_\omega(\Delta)=[\cos(\omega\Delta)\ \ \sin(\omega\Delta)]$, Gram matrices
$S_\omega=\mathbb E[x_\omega^\top x_\omega]$ and
$H_{\omega\nu}=\mathbb E[x_\omega^\top x_\nu]$ over a separation prior, and the
whitened cross-Gram $Q_{\omega\nu}=S_\omega^{-1/2}H_{\omega\nu}S_\nu^{-1/2}$,
the singular values of $Q_{\omega\nu}$ are the canonical correlations between
$V_\omega$ and $V_\nu$, and
\begin{equation}
c_{\omega\nu}=\tfrac12\lVert Q_{\omega\nu}\rVert_F^2\in[0,1]
\label{eq:collision}
\end{equation}
is invariant to phase rotation inside a pair and to any change of basis of
either subspace. For $\Delta\sim\mathrm{Unif}[0,L]$, $H_{\omega\nu}$ has a closed
form in $\sin t/t$ and $(1{-}\cos t)/t$ at $t=(\omega{\mp}\nu)L$
(App.~\ref{sec:proofs}).

\subsection{The spectral budget is a dimension count}
\label{sec:budget}

Block-whitening every pair gives a global correlation Gram $R$ with diagonal
blocks $I_2$. Write $\bar c$ for the mean of the $c_{ij}$.

\begin{theorem}[Spectral budget identity]
\label{thm:budget}
$\operatorname{tr}R = 2K$ and $\operatorname{tr}(R^2)=2K[1+(K{-}1)\bar c]$, so the
Rényi-$2$ (stable) effective dimension of the positional basis is exactly
\begin{equation}
r_2(R)=\frac{(\operatorname{tr}R)^2}{\operatorname{tr}(R^2)}
      =\frac{2K}{1+(K-1)\bar c}.
\label{eq:budget-identity}
\end{equation}
\end{theorem}

The nominal budget $2K$ is deflated by the average pairwise redundancy, and by
nothing else. This is the sense in which RoPE has a spectral budget, and it
names the variable a schedule acts on: \emph{once $K$ and the sampled endpoints
are fixed, allocation moves exactly one quantity, $\bar c$}. The exact-range
experiment of \S\ref{sec:exp-identify} manipulates precisely that. Note also
what \eqref{eq:budget-identity} does not do: it is symmetric in the frequencies
and cannot single out a particular density shape --- a prediction
\S\ref{sec:exp-tau} confirms.

The deflation is concentrated at the slow end.

\begin{proposition}[Low-frequency collapse]
\label{prop:collapse}
As $\omega L\to0$, $V_\omega\to\operatorname{span}\{1,\Delta\}$, and for slow
channels $x{=}\omega L$, $y{=}\nu L$,
$2-\lVert Q_{x,y}\rVert_F^2=\tfrac{19}{12600}(x^2{-}y^2)^2+O(\epsilon^6)$: slow
subspaces approach a common two-dimensional limit at fourth order. Under the
softmax metric $F=\operatorname{diag}(p)-pp^\top$, which annihilates the constant
direction, the centred limit is
$\operatorname{span}\{\Delta-\mathbb E_p\Delta,\ \Delta^2-\mathbb E_p\Delta^2\}$.
\end{proposition}

Concretely, at $L{=}4096$, $b{=}5{\times}10^5$, $K{=}64$ the $24$ pairs with
$\omega L\le1$ occupy $48$ nominal dimensions but have block-whitened
$r_2 = 2.00$: about $96\%$ of that sub-budget is redundant. Slow channels are
\emph{redundant}, not unused --- they still carry content, and
Prop.~\ref{prop:collapse} says only that many of them describe nearly the same
positional direction.

\subsection{The table is a training-time coordinate system}
\label{sec:coordinate}

Theorem~\ref{thm:budget} constrains what a table can express. It does not
predict what a trained model does with one, and the gap is not small. In a
frozen $50$M diagnostic ($L{=}512$, $b{=}5{\times}10^5$; no training, no
gradient) we cross two trained weight sets with two runtime tables
(Table~\ref{tab:coadapt}). Substituting \evq{} for the geometric table at
inference \emph{raises} the block-whitened stable rank from $4.57$ to $12.54$
while perplexity degrades from $7.14$ to $76.20$.

\begin{table}[tb]
\centering
\small
\setlength{\tabcolsep}{6pt}
\renewcommand{\arraystretch}{0.95}
\caption{\textbf{Static geometry does not predict trained behaviour.} Two
seed-$42$ $50$M models trained with the geometric and \evq{} tables, each
evaluated under both runtime tables with all parameters frozen ($1{,}920$
head--query observations). The worst cell is also the one whose static
$r_2$ improves most.}
\label{tab:coadapt}
\begin{tabular}{@{}llcc@{}}
\toprule
Trained weights & Runtime table & PPL & Rényi-$2$ rank $r_2$ \\
\midrule
Geo   & Geo   & $\mathbf{7.14}$ & $4.57$ \\
Geo   & \evq{} & $76.20$ & $\mathbf{12.54}$ \\
\evq{} & Geo   & $23.05$ & $4.57$ \\
\evq{} & \evq{} & $\mathbf{7.16}$ & $12.54$ \\
\bottomrule
\end{tabular}
\end{table}


Two consequences discipline the rest of the paper. First, no static criterion
--- collision, stable rank, or log-determinant --- may serve as a proxy for
trained quality, so every task claim in \S\ref{sec:experiments} is made on
trained models. Second, the two \emph{self-consistent} systems are essentially
tied in-window ($7.14$ vs.\ $7.16$) while both mismatched cells collapse: on LM
loss the table and weight main effects are $+0.599$ and $-0.597$, and their
interaction is $-3.537$, about $5.9\times$ larger (configuration-level bootstrap
CI $[-5.17,-3.04]$). The table behaves as a coordinate system fixed before
training, with the weights learning coefficients in it.

This is not a tuning failure that a better adapter would remove.

\begin{theorem}[Post-hoc transplant obstruction]
\label{thm:obstruction}
Let $R_\Omega(\Delta)$ be the block-rotation operator of a frequency multiset
$\Omega$. If position-independent invertible maps $A,B$ satisfy
$A^\top R_{\Omega'}(\Delta)B = R_\Omega(\Delta)$ for all $\Delta$ in a set
containing an open interval, then $\Omega'$ and $\Omega$ have the same frequency
multiset, up to sign, permutation, and change of basis inside a pair.
\end{theorem}

The proof is short: $\Delta{=}0$ gives $A^\top B = I$, so the relation becomes a
similarity $A^\top R_{\Omega'}A^{-\top}=R_\Omega$; differentiating at zero makes
the generators similar, and similarity preserves the spectrum $\{\pm i\omega_k\}$
(App.~\ref{sec:proofs}). A nontrivial table change therefore cannot be absorbed
exactly by any fixed linear Q/K compensation, and the in-window cost of
retrofitting a pretrained checkpoint is structural rather than an artefact of
insufficient adaptation. The theorem does not rule out approximate repair by
retraining, nonlinear maps, or the V/O and residual paths --- which is what the
adaptation protocols of \S\ref{sec:exp-mature} do.

\subsection{A closed-form allocation}
\label{sec:construction}

A positive density $\rho\in C^2([0,1])$ with $\int_0^1\rho=1$ induces a discrete
allocation by inverse-CDF quantisation, $\phi_k=F_\rho^{-1}(u_k)$; geometric RoPE
is $\rho\equiv1$. Choosing $\rho$ needs a tractable proxy for $\bar c$. The exact
kernel is oscillatory and is not pointwise approximated by a smooth two-term
model, so we analyse the convex surrogate
\begin{equation}
\Capp[\rho]=\frac{\alpha}{2}\int_0^1\!\rho^2\,d\phi
 +\frac{\beta}{2}\iint_{[0,1]^2}\!\!\rho(\phi)\rho(\psi)\min(\phi,\psi)\,d\phi\,d\psi,
\label{eq:Capp}
\end{equation}
with $\alpha,\beta>0$ fitted to the exact kernel on a finite channel grid
(App.~\ref{sec:tau-scaling}). The diagonal term penalises channel-load
concentration; $\min(\phi,\psi)$ is the Green kernel of $-\partial_\phi^2$ and
penalises pairs of channels that are \emph{both} slow --- the regime
Prop.~\ref{prop:collapse} identifies as redundant.

\begin{theorem}[Stationary allocation under the surrogate]
\label{thm:ode}
Under \eqref{eq:Capp} with $\alpha{>}0,\beta{\geq}0$ and constraints
$\rho\in C^2([0,1])$, $\rho{>}0$, $\int_0^1\rho{=}1$, the unique constrained
minimiser is
\begin{equation}
\rho_\tau(\phi)=\frac{\tau\cosh(\tau(1{-}\phi))}{\sinh\tau},
\qquad \tau=\sqrt{\beta/\alpha};
\label{eq:rho-tau}
\end{equation}
at $\beta{=}0$ it is the $\tau{\to}0$ limit $\rho_0\equiv1$.
\end{theorem}

\begin{corollary}[Closed-form quantiles]
\label{cor:quantiles}
$F_\tau(\phi)=1-\sinh(\tau(1{-}\phi))/\sinh\tau$, so
\begin{equation}
\phi_k(\tau)=1-\tfrac{1}{\tau}\operatorname{arcsinh}\!\left((1-u_k)\sinh\tau\right),
\qquad u_k=\tfrac{k+1/2}{K},
\label{eq:warp}
\end{equation}
which tends to $u_k$ as $\tau\to0$, recovering the uniform-log midpoint grid.
\end{corollary}

The full stationary density also admits a Fisher-forced particular solution; we
deploy the homogeneous \emph{pure-tether} branch \eqref{eq:rho-tau} because it
alone has an analytically invertible CDF (App.~\ref{subsec:fisher-ext}).

\paragraph{The recipe, and the status of $\tau$.} With $\deff$ the effective head
dimension, equal to $\dhd$ for the full-rotary architectures studied here,
\begin{equation}
\label{eq:evq-practical}
\tau = \deff/\sqrt{\Ltr},\qquad
u_k = (k+\tfrac12)/K,\qquad
\omega_k^{\mathrm{EVQ}} = b^{-\phi_k(\tau)} .
\end{equation}
Set $\tau$, evaluate \eqref{eq:warp}, replace the inverse frequencies: no
auxiliary loss, no extra parameters, no declared deployment length. A
post-softmax transport model supplies the $\deff/\sqrt{\Ltr}$ form but not its
constant, and its small-$\tau$ expansion does not certify the deployed
$\tau\!\approx\!4$ regime (App.~\ref{sec:leff-derivation}); we therefore treat
$c{=}1$ as a zero-search default, and across the $21$ configurations of
\S\ref{sec:exp-tau} the selected $\tau$ never leaves a $0.75\times$--$1.5\times$
window around it. Finally, two grids appear below and are \emph{not}
interchangeable: \evq{} deploys the midpoint grid of \eqref{eq:warp}, while the
endpoint-normalised variant $u_k{=}k/K$ ($\phi_0{=}0$ preserves $\omega_0{=}1$)
is a shape diagnostic used where a control must pin the sampled extrema exactly,
as in \S\ref{sec:exp-identify}.

\subsection{Where the finite budget moves}
\label{sec:budget-shift}

In wavelength $\lambda(\phi)=2\pi b^\phi$, a training span $L$ and evaluation
ceiling $M{>}L$ split the axis into three regions: $\lambda\le L$ completes at
least one cycle during training; $L<\lambda\le M$ varies across the target window
without completing one; $\lambda>M$ completes none even at the ceiling.

\begin{lemma}[Single-crossing budget shift]
\label{lem:budget-crossing}
For every $\tau>0$, $\rho_\tau$ crosses the uniform density exactly once, at
$\phi_c(\tau)=1-\tau^{-1}\operatorname{arcosh}(\sinh\tau/\tau)\le1-1/\sqrt3\approx0.4226$.
\end{lemma}

For OLMo-2 ($b{=}5{\times}10^5$, $L{=}4096$, $M{=}32768$) the boundaries
$\phi_L{=}0.494$ and $\phi_M{=}0.652$ lie to the right of that universal bound,
so at $\tau{=}2$ the $64$-pair table moves from $32/10/22$ channels to
$43/8/13$: eleven channels bought for the completed-cycle region by spending
density in both longer-wavelength regions. This is an exact budget statement,
not a trained-loss theorem, and it yields the falsifiable prediction tested in
\S\ref{sec:discussion}. Proof and a divergence inequality are in
App.~\ref{sec:waterbed-proof}.
